%! Comparing Studies and Choosing a Method — Unit 1, Lesson 1.7 — Exit Ticket (Answer Key)
\documentclass[10pt]{article}
\usepackage{saar-article}
\usepackage{saar-key}   % pulls in -boxes; do NOT also load -boxes
\newcommand{\TallMath}[1]{$\displaystyle #1\rule[-1.4em]{0pt}{3.2em}$}

\begin{document}

\pageheader{Unit 1, Lesson 1.7}{Exit Ticket --- Answer Key}
% NAMESTRIP: no \namedateperiod here — mirrors the blank. Teacher-only prose goes
% in the lesson plan as \begin{teachernote}[Exit Ticket], never in a key.

\vspace{0.15in}

\begin{notesbox}{Exit Ticket --- No Notes}
\small
Bayside Middle School has two studies on file. \textbf{Study A:} the counselor
compared the $150$ students in band with the $600$ who are not --- $90$ of the
band students and $240$ of the others had a B average or better. \textbf{Study
B:} $80$ sixth-graders volunteered and a random number generator sent $40$ to a
$10$-minute silent reading block and $40$ to their usual homeroom; $70\%$ of the
reading-block group improved their reading score against $40\%$ of the homeroom
group.

\begin{enumerate}[leftmargin=1.6em, itemsep=8pt, topsep=3pt]

  \item Sort the two studies.

        \textbf{(a)} What percent of the band students had a B average or
        better?

        \begin{work}
          \dfrac{90}{150} &= 0.60 \\
                          &= 60\%
        \end{work}

        \textbf{(b)} Study A is a(n) \ans{observational study} \hfill Study B is
        a(n) \ans{experiment}

        \textbf{(c)} Only one of these studies may say the program
        \emph{caused} the difference. Which one, and what earns it that word?
        \ansline{Study B. A generator decided who got the reading block, so the two groups}
        \ansline{started out alike and the block is the only thing left to explain the gap.}

  \item The principal asks a new question: \emph{do students who ride the bus
        more than $45$ minutes each way score lower on the state test?}

        \textbf{(a)} Could she assign students to a long bus ride? \ans{No}

        \textbf{(b)} Say why or why not, and name the study type she must run
        instead.
        \ansline{She cannot assign where a family lives --- so it must be an observational study.}

  \item She also wants to know \emph{why} sixth-graders skip breakfast, in their
        own words. Name the collection method that fits best, and say why a
        written survey given to all $750$ students is a worse fit for this
        particular question.
        \ansline{Interviews (a focus group also works). ``Why'' needs depth and follow-up questions.}
        \ansline{A survey gives everyone the same fixed boxes and cannot ask what a student means.}

\end{enumerate}
\end{notesbox}

\end{document}
